%Version 3.1 December 2024
% See section 11 of the User Manual for version history
%
%%%%%%%%%%%%%%%%%%%%%%%%%%%%%%%%%%%%%%%%%%%%%%%%%%%%%%%%%%%%%%%%%%%%%%
%%                                                                 %%
%% Please do not use \input{...} to include other tex files.       %%
%% Submit your LaTeX manuscript as one .tex document.              %%
%%                                                                 %%
%% All additional figures and files should be attached             %%
%% separately and not embedded in the \TeX\ document itself.       %%
%%                                                                 %%
%%%%%%%%%%%%%%%%%%%%%%%%%%%%%%%%%%%%%%%%%%%%%%%%%%%%%%%%%%%%%%%%%%%%%

%%\documentclass[referee,sn-basic]{sn-jnl}% referee option is meant for double line spacing

%%=======================================================%%
%% to print line numbers in the margin use lineno option %%
%%=======================================================%%

%%\documentclass[lineno,pdflatex,sn-basic]{sn-jnl}% Basic Springer Nature Reference Style/Chemistry Reference Style

%%=========================================================================================%%
%% the documentclass is set to pdflatex as default. You can delete it if not appropriate.  %%
%%=========================================================================================%%

%%\documentclass[sn-basic]{sn-jnl}% Basic Springer Nature Reference Style/Chemistry Reference Style

%%Note: the following reference styles support Namedate and Numbered referencing. By default the style follows the most common style. To switch between the options you can add or remove “Numbered” in the optional parenthesis. 
%%The option is available for: sn-basic.bst, sn-chicago.bst%  
 
%%\documentclass[pdflatex,sn-nature]{sn-jnl}% Style for submissions to Nature Portfolio journals
%%\documentclass[pdflatex,sn-basic]{sn-jnl}% Basic Springer Nature Reference Style/Chemistry Reference Style
\documentclass[pdflatex,sn-mathphys-num]{sn-jnl}% Math and Physical Sciences Numbered Reference Style
%%\documentclass[pdflatex,sn-mathphys-ay]{sn-jnl}% Math and Physical Sciences Author Year Reference Style
%%\documentclass[pdflatex,sn-aps]{sn-jnl}% American Physical Society (APS) Reference Style
%%\documentclass[pdflatex,sn-vancouver-num]{sn-jnl}% Vancouver Numbered Reference Style
%%\documentclass[pdflatex,sn-vancouver-ay]{sn-jnl}% Vancouver Author Year Reference Style
%%\documentclass[pdflatex,sn-apa]{sn-jnl}% APA Reference Style
%%\documentclass[pdflatex,sn-chicago]{sn-jnl}% Chicago-based Humanities Reference Style

%%%% Standard Packages
%%<additional latex packages if required can be included here>

\usepackage{graphicx}%
\usepackage{multirow}%
\usepackage{amsmath,amssymb,amsfonts}%
\usepackage{amsthm}%
\usepackage{mathrsfs}%
\usepackage[title]{appendix}%
\usepackage{xcolor}%
\usepackage{textcomp}%
\usepackage{manyfoot}%
\usepackage{booktabs}%
\usepackage{algorithm}%
\usepackage{algorithmicx}%
\usepackage{algpseudocode}%
\usepackage{listings}%
\usepackage{subcaption}
\usepackage{adjustbox}
\usepackage{multicol,multirow}
%%%%

%%%%%=============================================================================%%%%
%%%%  Remarks: This template is provided to aid authors with the preparation
%%%%  of original research articles intended for submission to journals published 
%%%%  by Springer Nature. The guidance has been prepared in partnership with 
%%%%  production teams to conform to Springer Nature technical requirements. 
%%%%  Editorial and presentation requirements differ among journal portfolios and 
%%%%  research disciplines. You may find sections in this template are irrelevant 
%%%%  to your work and are empowered to omit any such section if allowed by the 
%%%%  journal you intend to submit to. The submission guidelines and policies 
%%%%  of the journal take precedence. A detailed User Manual is available in the 
%%%%  template package for technical guidance.
%%%%%=============================================================================%%%%

%% as per the requirement new theorem styles can be included as shown below
\theoremstyle{thmstyleone}%
\newtheorem{theorem}{Theorem}%  meant for continuous numbers
%%\newtheorem{theorem}{Theorem}[section]% meant for sectionwise numbers
%% optional argument [theorem] produces theorem numbering sequence instead of independent numbers for Proposition
\newtheorem{proposition}[theorem]{Proposition}% 
%%\newtheorem{proposition}{Proposition}% to get separate numbers for theorem and proposition etc.

\theoremstyle{thmstyletwo}%
\newtheorem{example}{Example}%
\newtheorem{remark}{Remark}%

\theoremstyle{thmstylethree}%
\newtheorem{definition}{Definition}%

\raggedbottom
%%\unnumbered% uncomment this for unnumbered level heads

\begin{document}

\title[Robust Nighttime Vehicle Localization Under Headlight Saturation via Ensemble Detection and Adaptive 3-D Tracking]{Robust Nighttime Vehicle Localization Under Headlight Saturation via Ensemble Detection and Adaptive 3-D Tracking}

%%=============================================================%%
%% GivenName	-> \fnm{Joergen W.}
%% Particle	-> \spfx{van der} -> surname prefix
%% FamilyName	-> \sur{Ploeg}
%% Suffix	-> \sfx{IV}
%% \author*[1,2]{\fnm{Joergen W.} \spfx{van der} \sur{Ploeg} 
%%  \sfx{IV}}\email{iauthor@gmail.com}
%%=============================================================%%

\author[1]{\fnm{Nguyen Van} \sur{Huan}}
\email{huan.nv2416500@sis.hust.edu.vn}
\author[1]{Duong Tan Minh}
\author[1]{Vu Van Giang}
\author[1]{Nguyen Thanh Kien}
\author[1]{Dang Quang Huy}
\author*[1]{Ho Viet Duc Luong}
\email{luonghvd@soict.hust.edu.vn}

\affil[1]{\orgdiv{School of Information and Communications Technology}, \orgname{Hanoi University of Science and Technology}, \orgaddress{\street{1 Dai Co Viet Street}, \city{Hanoi}, \postcode{100000}, \country{Vietnam}}}


%%==================================%%
%% Sample for unstructured abstract %%
%%==================================%%

\abstract{The abstract serves both as a general introduction to the topic and as a brief, non-technical summary of the main results and their implications. Authors are advised to check the author instructions for the journal they are submitting to for word limits and if structural elements like subheadings, citations, or equations are permitted.}

%%================================%%
%% Sample for structured abstract %%
%%================================%%

% \abstract{\textbf{Purpose:} The abstract serves both as a general introduction to the topic and as a brief, non-technical summary of the main results and their implications. The abstract must not include subheadings (unless expressly permitted in the journal's Instructions to Authors), equations or citations. As a guide the abstract should not exceed 200 words. Most journals do not set a hard limit however authors are advised to check the author instructions for the journal they are submitting to.
% 
% \textbf{Methods:} The abstract serves both as a general introduction to the topic and as a brief, non-technical summary of the main results and their implications. The abstract must not include subheadings (unless expressly permitted in the journal's Instructions to Authors), equations or citations. As a guide the abstract should not exceed 200 words. Most journals do not set a hard limit however authors are advised to check the author instructions for the journal they are submitting to.
% 
% \textbf{Results:} The abstract serves both as a general introduction to the topic and as a brief, non-technical summary of the main results and their implications. The abstract must not include subheadings (unless expressly permitted in the journal's Instructions to Authors), equations or citations. As a guide the abstract should not exceed 200 words. Most journals do not set a hard limit however authors are advised to check the author instructions for the journal they are submitting to.
% 
% \textbf{Conclusion:} The abstract serves both as a general introduction to the topic and as a brief, non-technical summary of the main results and their implications. The abstract must not include subheadings (unless expressly permitted in the journal's Instructions to Authors), equations or citations. As a guide the abstract should not exceed 200 words. Most journals do not set a hard limit however authors are advised to check the author instructions for the journal they are submitting to.}

\keywords{keyword1, Keyword2, Keyword3, Keyword4}

%%\pacs[JEL Classification]{D8, H51}

%%\pacs[MSC Classification]{35A01, 65L10, 65L12, 65L20, 65L70}

\maketitle

\section{Introduction}
\label{sec:intro}
%% ==============================================================

Nighttime driving remains a disproportionately hazardous operating
condition for road users. Although nighttime traffic volume is
substantially lower than daytime traffic volume, severe crashes occur
more frequently at night, and glare from oncoming headlights is widely
recognized as a major contributor to degraded visual
performance~\cite{fhwa2025nightvisibility,hwang2018headlightglare,beard2025glarelitreview}.
This challenge has motivated the development of intelligent front-lighting systems that seek to preserve the host driver's visibility while
reducing glare exposure for other road users. In particular, \emph{Adaptive Driving Beam} (ADB) systems selectively suppress only the portion of the high beam that would dazzle an oncoming vehicle, thereby retaining as much useful illumination as possible outside the affected region~\cite{Neumann2014ADB,Shin2026DOE_AFS}.

From a computer-vision perspective, however, ADB is not merely a lighting problem. It is a closed-loop perception-and-actuation task in which the system must first localize the relevant oncoming road user in the image, then translate that localization into a beam-suppression pattern on a discrete headlight actuator. This coupling introduces a practical failure mode that remains insufficiently addressed in the literature: when the front-facing camera is exposed to intense oncoming headlights, blooming, flare, specular reflection, and rapid auto-exposure adaptation can severely degrade the visual signal. In the
most adverse cases, the detector output becomes unstable or vanishes entirely, causing the suppression region to flicker or disappear. As a
result, a camera-centric ADB pipeline may fail precisely in the regime where robust localization is most safety-critical.

This limitation highlights a broader gap between mainstream object detection and the requirements of perception-guided beam control. Recent nighttime vehicle detection studies have shown that low-light driving scenes differ fundamentally from daytime imagery: salient visual evidence is often concentrated in sparse, high-intensity headlight regions rather than in stable object texture or contour structure.
Accordingly, recent approaches have explored specialized feature design, headlight-aware modeling, and highlight-guided attention to improve nighttime vehicle detection under challenging illumination conditions \cite{zhang2024vision,lashkov2025machine,song2025highlightnet}.
At the same time, the automotive perception literature has increasingly
demonstrated the value of radar--camera fusion for maintaining robust
localization when vision alone becomes unreliable, particularly under
adverse weather, darkness, and low-visibility
conditions~\cite{yao2023radar,deng2025advances,xiao2025radar}.
Yet, these two strands of work are rarely brought together in the specific context of ADB, where the goal is not a general scene understanding, but a stable continuation of the suppression zone during brief periods of visual collapse. A second limitation of existing work is methodological. Most vision papers evaluate nighttime detection in pixel space using standard metrics such as Precision, Recall, and mAP, whereas automotive lighting studies emphasize photometric or regulatory performance. For perception-driven matrix headlight control, neither perspective alone is
fully sufficient. Small image-space localization errors may map to incorrect LED cells, producing either residual glare or unnecessary darkening. Moreover, image-level detection scores do not directly capture whether a downstream suppression actuator behaves correctly during transient detector dropout. This motivates the need for an
evaluation protocol that links perception robustness to actuation fidelity in a reproducible closed-loop setting.

To address these gaps, we propose  \textbf{AuraBeam}, an end-to-end nighttime vehicle localization and glare-suppression framework designed for camera-centric ADB under headlight saturation. Rather than treating
nighttime detection as an isolated image-space task, AuraBeam combines detection, adaptive tracking, and physical actuation into a unified pipeline. First, an ensemble vision module runs YOLOv5 and YOLOv8 in parallel and fuses their predictions using Weighted Non-Maximum
Suppression, improving robustness under blooming and adverse illumination. Second, a six-state 3-D Kalman filter integrates the
image-space observations with a physics-based depth surrogate (Pseudo-Radar) and switches to a radar-only observation mode during
camera degradation, allowing the suppression trajectory to be propagated through short intervals of visual dropout. Third, the framework is validated in a compact Hardware-in-the-Loop (HIL) environment and evaluated using a new actuation-level metric, the \emph{Glare Suppression Success Rate} (GSSR), which measures suppression fidelity on the LED grid rather than only in pixel space.

The main contributions of this work are summarized as follows:
\begin{itemize}
    \item \textbf{Glare-robust nighttime localization via ensemble detection.}
    We design a dual-model YOLO pipeline in which YOLOv5 and YOLOv8 are executed in parallel and fused by Weighted Non-Maximum Suppression. This design improves robustness to blooming, reflection, and
    illumination instability, leading to higher Recall in challenging nighttime scenarios without requiring a substantially more complex downstream control pipeline.

    \item \textbf{Adaptive 3-D tracking under temporary visual collapse.}
    We introduce a six-state 3-D Kalman filtering framework that fuses image-space detections with a physics-based depth surrogate and adaptively switches its observation model when the camera signal
    becomes unreliable. This enables stable continuation of the suppression region during short periods of detection dropout and reduces localization jitter in adverse nighttime conditions.

    \item \textbf{Actuation-aware evaluation in a closed-loop HIL setting.}
    We formulate the \emph{Glare Suppression Success Rate (GSSR)} as an LED-grid-based metric for evaluating downstream suppression quality, and we construct a reproducible HIL evaluation framework with a
    manually annotated Golden Test Set spanning five adversarial scenarios, including curved roads, rain, fog/snow, and abrupt
    exposure disruption.
\end{itemize}

Overall, this work aims to bridge three literatures that are often studied separately: nighttime vehicle detection, robust tracking/fusion under degraded sensing, and intelligent headlight actuation. By
focusing on the specific failure mode of \emph{glare-induced detection collapse}, we show that improving camera-centric ADB requires not only better detection, but also graceful localization continuity and actuation-aware evaluation.
The remainder of this paper is organized as follows.
Section~\ref{sec:related} reviews the relevant literature and positions the proposed work within recent advances in nighttime vehicle detection, sensor fusion, and intelligent headlight systems.
Section~\ref{sec:method} details the proposed AuraBeam framework.
Section~\ref{sec:eval} presents the evaluation protocol, datasets, and performance metrics.
Section~\ref{sec:results} reports and analyzes the experimental results.
Finally, Section~\ref{sec:conclusion} summarizes the main findings and
outlines directions for future work.

\section{Background and Related Work}
\label{sec:related}
%% ==============================================================

% \subsection{Adaptive Driving Beam Systems}
% \label{subsec:adb}

Adaptive and intelligent headlight systems have evolved from simple Auto High Beam (AHB) switching toward spatially selective beam shaping that preserves forward visibility while reducing glare to other road users. A recent review by Nkrumah \emph{et al.} identifies camera-based scene understanding, rule-based beam modulation, and adaptive intensity
control as the dominant paradigm in intelligent headlight systems~\cite{nkrumah2024review}. In parallel, recent optical and regulatory studies have emphasized that production ADB/AFS designs must satisfy not only visibility requirements but also strict glare-compliance constraints under ECE-R123-like frameworks~\cite{shin2026optimization}. These developments confirm that modern headlight control is no longer a purely mechanical or optical problem; it is increasingly a perception-and-control problem in which the quality of the front-facing vision stack directly determines beam-selection safety.
Despite this progress, most deployed and prototype ADB pipelines remain
fundamentally \emph{camera-centric}: an oncoming vehicle is first detected
in the image plane, after which a suppression mask or beam cut-off region is generated for a segmented LED or pixelated headlamp. This architecture works well under moderate nighttime conditions, but becomes fragile when the camera is saturated by strong oncoming headlights, specular reflections, or abrupt exposure changes. A recent UK literature review on vehicle-lighting glare further reinforces that glare remains an unresolved practical safety issue influenced by light source intensity, headlamp aim, road geometry, vehicle height, and retrofitted lighting
configurations~\cite{beard2025glarelitreview}. Thus, while recent literature has improved the design of adaptive beam systems, a key systems-level gap remains: how to maintain glare suppression when the vision channel itself
temporarily collapses.

% \subsection{Nighttime Vehicle Detection Under Blooming and Highlight Saturation}
% \label{subsec:detection}

Nighttime vehicle detection is substantially harder than daytime detection because the scene is dominated by sparse, high-intensity light
sources rather than stable object texture. In such conditions, blooming, flare, reflection, and exposure adaptation distort the apparent spatial extent of the headlights and often erase the geometric contours of the vehicle body. Recent work has therefore moved beyond generic daytime detectors toward methods that explicitly exploit nighttime-specific
appearance cues. For example, Zhang \emph{et al.} proposed an improved HOG-based framework for on-road nighttime vehicle detection and tracking, showing that handcrafted structural cues can still be effective when illumination is poor and vehicle contours are weak~\cite{zhang2024vision}.
More recently, Lashkov \emph{et al.} developed a learning-based pipeline that couples headlight extraction with low-illumination tracking, again highlighting that vehicle lights themselves become the dominant visual
carrier in dark scenes~\cite{lashkov2025machine}.
Deep models have also begun to incorporate this observation explicitly. HighlightNet, introduced by Song \emph{et al.} \cite{song2025highlightnet}, formulates nighttime vehicle detection as a dual-branch problem that jointly reasons about vehicle presence and highlight regions, using highlight-guided attention to force the detector to focus on the most informative luminous structures. These recent studies collectively suggest an important shift in the literature: nighttime detection should not be treated as a simple low-light version of daytime detection, but as a distinct visual regime in which the detector must reason under partial loss of texture, extreme local saturation, and ambiguous light-source morphology.
However, most of these methods still evaluate success primarily in pixel-space, using standard detection metrics such as Precision, Recall, or mAP. For ADB control, this is only part of the problem. A detector may remain acceptable in pixel-space while still producing unstable or misaligned suppression decisions once mapped onto a discrete LED grid. Moreover, many nighttime detection papers focus on \emph{detecting} vehicles, whereas ADB requires \emph{persistent localization} across short intervals of visual degradation. This distinction motivates our use of a
dual-model detector and an actuation-aware downstream evaluation, rather
than relying on detection performance alone.

% \subsection{Trajectory Tracking and Radar--Camera-Inspired Fusion}
% \label{subsec:kf}

Tracking plays a central role whenever detections are intermittent, noisy, or partially missing. Classical state estimators such as the Kalman Filter (KF) remain attractive because they provide a lightweight and an interpretable mechanism for propagating object state through short periods of observation loss. Yet, in automotive perception, the more important recent trend is not simply filtering, but \emph{multimodal
fusion}. A comprehensive review by Yao \emph{et al.} shows that radar--camera fusion has become a major paradigm for autonomous-driving perception because cameras provide rich semantics while radar offers
depth and motion cues that are far more robust to darkness, adverse weather, and poor illumination~\cite{yao2023radar}. Recent
surveys likewise emphasize that multimodal fusion improves robustness precisely in the regimes where camera-only perception is most fragile \cite{qian2025review}.
Recent fusion methods span feature-level, decision-level, and geometric fusion designs. Deng \emph{et al.} review recent radar--camera object detection advances for autonomous driving and note that the two modalities are complementary in both sensing physics and failure characteristics~\cite{deng2025advances}. Xiao \emph{et al.} \cite{xiao2025radar} further show that combining perspective-view and bird's-eye-view cues improves 3-D object detection performance, highlighting that depth-aware fusion is increasingly central to robust vehicle localization. Even when full radar hardware is not available, these studies still motivate a key design principle: \emph{localization should not depend exclusively on the optical channel}
when the target application is safety-critical and must operate under glare, rain, or low-light perturbations.
At the same time, the specific fusion problem in ADB is more constrained than general autonomous driving. The objective is not full-scene 3-D understanding, but reliable continuation of the suppression region for
the most relevant oncoming light source over short time horizons. This creates two unique requirements. First, when the camera becomes saturated,
the image-plane coordinates become weakly observed or unobservable. Second, any auxiliary depth signal must fail independently from the camera; otherwise, both modalities collapse under the same glare event. Our design follows the logic of radar--camera fusion literature, but in a cost-constrained form: we combine image-space detections with a physics-based depth surrogate and switch to a radar-only observation model during visual dropout. In this sense, our method is not intended to replace full radar--camera fusion, but to instantiate its robustness principle in an ADB-specific and low-cost setting.

% \subsection{Hardware-in-the-Loop and Vision-System Validation}
% \label{subsec:hil}

A second gap in the literature concerns \emph{how} intelligent headlight systems are evaluated. Traditional automotive lighting studies often rely
on photometric bench tests, regulatory compliance analysis, or expensive on-road experiments, whereas computer-vision papers typically stop at image-space metrics. For systems such as ADB, neither perspective is fully sufficient: the former underrepresents perception failure modes, while the latter may ignore the downstream consequences of localization errors on actual beam actuation.
Recent validation literature increasingly advocates testbench-based hardware-in-the-loop (HIL), software-in-the-loop (SIL), and vehicle-in-the-loop (VIL) methodologies for ADAS and autonomous-vehicle
development \cite{cheng2024survey, langer2026vehicle,balan2025perspective}. Cheng \emph{et al.} survey VIL simulation testing for autonomous vehicles and emphasize its role in bridging pure simulation and costly full-road experiments through repeatable, physically grounded
testing setups~\cite{cheng2024survey}. More recently, Langer \emph{et al.} demonstrate a VIL framework specifically for front-camera verification using adaptive high beam, underscoring the growing need for
reproducible evaluation of perception-dependent lighting functions~\cite{langer2026vehicle}. Broader recent perspectives on SIL/HIL in automotive testing also argued that closed-loop virtual--physical validation is increasingly necessary as perception, control, and embedded implementation become tightly coupled~\cite{balan2025perspective}.

These developments strongly support the use of a compact HIL platform in our setting. However, existing VIL/HIL work on lighting is still primarily concerned with functional verification or platform-level testing, and much less with the specific computer-vision failure mode in which glare destroys the observation that is needed to maintain the beam
suppression mask. Our HIL framework is designed to address exactly this gap. By coupling the perception stack, the adaptive tracking logic, and a physical LED actuation layer, it enables evaluation at the level that ultimately matters for ADB control: whether the correct suppression zone is maintained in the presence of brief but safety-critical visual collapse. This is also why we later introduce an actuation-level metric
rather than relying exclusively on conventional detection scores.

% \subsection{Summary of the Literature Gap}
% \label{subsec:gap_summary}

In summary, recent literature has made clear progress in four directions: adaptive headlight control \cite{nkrumah2024review,shin2026optimization}, nighttime vehicle detection under challenging illumination conditions \cite{zhang2024vision,lashkov2025machine,song2025highlightnet}, radar--camera fusion for robust automotive
perception~\cite{yao2023radar,deng2025advances,xiao2025radar}, and HIL/VIL validation for intelligent vehicle
systems~\cite{cheng2024survey,langer2026vehicle,balan2025perspective}. Yet the intersection of these strands remains underexplored. In particular, the literature still lacks a low-cost, end-to-end framework that
addresses \emph{glare-induced detection collapse} in camera-centric ADB, maintains localization through short periods of observation loss, and evaluates the result at the level of discrete headlight actuation rather than only in pixel space. Our work is intended to fill this gap.

%% ==============================================================
\section{System Methodology}
\label{sec:method}
%% ==============================================================

\subsection{System Overview}
\label{subsec:overview}

Figure~\ref{fig:tong_quat} presents the overall AuraBeam pipeline. The system is organized as a closed-loop perception-to-actuation process comprising five modules: (A) an ensemble vision front-end for nighttime vehicle detection, (B) a physics-based depth surrogate referred to as Pseudo-Radar, (C) an adaptive 3-D Kalman filtering module for state estimation under partial observation loss, (D) a zone-mapping and anti-flicker controller that converts the estimated target state into a discrete suppression pattern, and (E) a hardware actuation layer based on an Arduino-driven LED matrix.

Unlike a conventional camera-only ADB pipeline, AuraBeam is designed to preserve localization continuity when the image stream becomes unreliable due to blooming, exposure disruption, or severe glare. The detector provides image-plane observations when confidence is sufficient, whereas the state estimator continues propagating the suppression trajectory during short intervals of visual collapse by relying on its internal motion model and the auxiliary depth signal. The output of this estimation stage is then translated into a suppression region on an $8\times8$ LED grid and sent to the hardware layer for real-time actuation.

\begin{figure*}[t]
    \centering
    \begin{subfigure}[b]{0.3\textwidth}
        \centering
        \includegraphics[width=\textwidth]{flow_charts/yolo_base.png}
        \caption{YOLO-only baseline pipeline.}
        \label{fig:anh_1}
    \end{subfigure}
    \hspace{2cm}
    \begin{subfigure}[b]{0.3\textwidth}
        \centering
        \includegraphics[width=\textwidth]{flow_charts/aurabeam.drawio.png}
        \caption{AuraBeam full system pipeline.}
        \label{fig:anh_2}
    \end{subfigure}
    \caption{Architectural comparison between the YOLO-only baseline pipeline and the proposed AuraBeam pipeline.}
    \label{fig:tong_quat}
\end{figure*}

To realize the Hardware-in-the-Loop (HIL) setup, we employ a low-cost peripheral hardware stack consisting of an Arduino Uno R3, an $8\times8$ LED matrix, a USB camera, and a serial communication interface. The processing unit executes the perception and filtering modules, while the microcontroller serves as the physical actuation layer. As summarized in Table~\ref{tab:bom}, the total peripheral cost is approximately \$23, excluding the host computer. This compact implementation is intended as a reproducible experimental platform rather than a production-grade automotive lighting system.

\begin{table}[t]
\centering
\caption{Hardware bill of materials.}
\label{tab:bom}
\begin{tabular}{llr}
\hline
\textbf{Component} & \textbf{Specification} & \textbf{Est. Cost (USD)} \\
\hline
Processing unit & PC or laptop with discrete GPU & --- \\
Microcontroller & Arduino Uno R3 & \$5 \\
LED matrix & 1588BS $8\times8$ module & \$1 \\
Camera & USB 1080p webcam or dashcam & \$15 \\
Serial interface & USB-A to USB-B cable & \$2 \\
\hline
\textbf{Total} & & \textbf{\$23} \\
\hline
\end{tabular}
\end{table}

For easier comprehension, Table \ref{tab:notations} summarizes the notations and symbols used throughout the papers.

\begin{table}[t]
\centering
\caption{Main symbols and notations.}
\label{tab:notations}
\begin{tabular}{ll}
\hline
\textbf{Symbol} & \textbf{Description} \\
\hline
$I_t$ & Input frame at time $t$ \\
$\mathcal{M}_1, \mathcal{M}_2$ & Detection models in the ensemble pipeline \\
$\mathcal{D}_i$ & Detection set produced by model $\mathcal{M}_i$ \\
$b_j=[c_x,c_y,w,h]^\top$ & Bounding box in center-width-height format \\
$s_j$ & Confidence score of a detection \\
$\hat{b}, \hat{s}$ & Fused bounding box and fused confidence score \\
$\mathbf{d}_t$ & Image-space observation vector at time $t$ \\
$\mathbf{x}_k$ & 3-D Kalman filter state vector at step $k$ \\
$\mathbf{F}$ & State transition matrix \\
$\mathbf{Q}$ & Process noise covariance matrix \\
$\mathbf{H}_{\mathrm{full}}$ & Observation matrix in full-observation mode \\
$\mathbf{H}_{\mathrm{occ}}$ & Observation matrix in occlusion-aware mode \\
$\mathbf{R}_{\mathrm{full}}$ & Measurement covariance in full-observation mode \\
$\mathbf{R}_{\mathrm{occ}}$ & Measurement covariance in occlusion-aware mode \\
$Z(t)$ & Pseudo-Radar depth prior at time $t$ \\
$\mathcal{B}_t$ & Suppression region on the LED grid at time $t$ \\
$\mathbf{G}_t$ & Output LED suppression grid at time $t$ \\
$\tau_{\mathrm{NMS}}$ & IoU threshold for Weighted NMS \\
$\tau_{\mathrm{conf}}$ & Confidence threshold for observation switching \\
$T_h$ & Hysteresis hold-time window \\
\hline
\end{tabular}
\end{table}
\clearpage
\subsection{Module A: Ensemble YOLO Vision Pipeline}
\label{subsec:yolo}

The first stage of AuraBeam performs nighttime vehicle detection using an ensemble of two complementary object detectors. Let $I_t$ denote the input frame at time $t$. Two models, denoted by $\mathcal{M}_1$ and $\mathcal{M}_2$, are executed in parallel on $I_t$. The first model is specialized for high-glare headlight patterns, whereas the second model is trained on a broader nighttime vehicle dataset to improve generalization to diverse scene structure and vehicle appearance.

Each model produces a set of detections $\mathcal{D}_i = \{(b_j, s_j)\}_{j=1}^{N_i}$, where $b_j = [c_x, c_y, w, h]^\top$ is a bounding box in center-width-height format and $s_j \in [0,1]$ is the associated confidence score. Because the two detectors are trained on purposefully different data distributions, their outputs are expected to be complementary, particularly under low-light conditions in which one detector may better preserve light-source sensitivity while the other better captures vehicle shape cues.

To combine the predictions, AuraBeam applies Weighted Non-Maximum Suppression (Weighted NMS). For two overlapping detections $(b_a, s_a)$ and $(b_b, s_b)$ that satisfy $\mathrm{IoU}(b_a, b_b) \geq \tau_{\mathrm{NMS}}$, the fused bounding box $\hat{b}$ and fused confidence $\hat{s}$ are computed as

\begin{equation}
\hat{b} = \frac{w_1 s_a b_a + w_2 s_b b_b}{w_1 s_a + w_2 s_b},
\label{eq:nms_box}
\end{equation}

\begin{equation}
\hat{s} = \frac{w_1 s_a + w_2 s_b}{w_1 + w_2},
\label{eq:nms_score}
\end{equation}
where $w_1, w_2 \in \mathbb{R}_{>0}$ are detector-specific weights selected on the validation set, and $\tau_{\mathrm{NMS}}$ is fixed to $0.5$. The output of Module A at frame $t$ is the image-space observation vector

\begin{equation}
\mathbf{d}_t = [c_x, c_y, w, h, \hat{s}]^\top \in \mathbb{R}^5.
\label{eq:det_vec}
\end{equation}

This representation retains both geometric information and a confidence signal that is later used by the adaptive observation logic in the Kalman filter.

\subsubsection*{Training Configuration}
\label{subsubsec:training}

The two detectors are trained independently in order to maximize feature diversity.

\paragraph{Model $\mathcal{M}_1$ (YOLOv5).}
Model $\mathcal{M}_1$ is trained on a custom manually annotated dataset of 5{,}500 images emphasizing high-glare headlight regions. Training is performed with an input resolution of $1024\times1024$, a batch size of 16, and a maximum of 100 epochs. Early stopping with patience 30 is used to reduce overfitting.

\paragraph{Model $\mathcal{M}_2$ (YOLOv8).}
Model $\mathcal{M}_2$ is fine-tuned on a second in-house dataset containing 6{,}000 manually annotated images. This model is trained for 100 epochs with an input resolution of 1024 and batch size 16. To improve robustness under nighttime imaging variability, the training pipeline includes Mosaic, MixUp, Copy-Paste, random scaling, HSV-based appearance perturbation, small-angle rotation, and weak perspective augmentation.

After training, the ensemble weights $w_1$ and $w_2$ in Eqs.~\eqref{eq:nms_box} and \eqref{eq:nms_score} are selected on the validation split to optimize downstream detection performance.

\subsection{Module B: Pseudo-Radar Depth Estimation}
\label{subsec:radar}

A key requirement of AuraBeam is to maintain target localization when the visual signal temporarily degrades. Instead of introducing a dedicated physical radar or LiDAR sensor, which would increase hardware cost and integration complexity, we construct a lightweight auxiliary depth signal termed \emph{Pseudo-Radar}. This signal is not intended to replace real range sensing; rather, it provides a simple depth prior that remains independent of instantaneous image brightness and can therefore continue to evolve during visual saturation.

Assuming that the oncoming vehicle follows an approximately constant approach velocity over short time intervals, the estimated range at time $t$ is defined as

\begin{equation}
Z(t) = Z_0 - v_{\mathrm{app}} t,
\label{eq:radar_z}
\end{equation}
where $Z_0$ is the initial range estimated at first acquisition and $v_{\mathrm{app}}$ is the relative approach velocity. Since $Z(t)$ is propagated analytically from the initial condition and elapsed time, it is unaffected by camera glare or short-term detector failure.

The approximation in Eq.~\eqref{eq:radar_z} is accurate only when the relative motion is close to constant velocity. If the true longitudinal acceleration is $a_{\mathrm{real}} \neq 0$, then the actual range evolves as

\begin{equation}
Z_{\mathrm{true}}(t) = Z_0 - v_{\mathrm{app}} t - \frac{1}{2} a_{\mathrm{real}} t^2,
\label{eq:radar_true}
\end{equation}
and the resulting depth estimation error is
\begin{equation}
\Delta Z(t) = Z(t) - Z_{\mathrm{true}}(t) = -\frac{1}{2} a_{\mathrm{real}} t^2.
\label{eq:depth_error}
\end{equation}

This error grows quadratically with time. For example, under moderate braking with $a \approx -0.5g$, the absolute depth error reaches approximately $2.45$ m after one second of occlusion, and under stronger braking with $a \approx -0.9g$, it may reach approximately $4.4$ m. Consequently, the Pseudo-Radar module should be interpreted as a low-cost temporal depth surrogate suitable for short dropouts, not as a substitute for production-grade radar sensing.

\subsection{Module C: Adaptive 3-D Kalman Filter}
\label{subsec:kf_fusion}

\subsubsection{State Representation and Motion Model}

To integrate image-plane localization with the auxiliary depth signal, AuraBeam maintains a six-dimensional state vector

\begin{equation}
\mathbf{x}_k = [c_x, \dot{c}_x, c_y, \dot{c}_y, Z, \dot{Z}]^\top \in \mathbb{R}^6,
\label{eq:state}
\end{equation}

where $(c_x, c_y)$ denote the image-space centroid of the target, $(\dot{c}_x, \dot{c}_y)$ are the corresponding image-plane velocities, and $(Z, \dot{Z})$ represent depth and longitudinal velocity. A constant-velocity state transition model is used:

\begin{equation}
\mathbf{F} =
\begin{bmatrix}
1 & \Delta t & 0 & 0 & 0 & 0 \\
0 & 1 & 0 & 0 & 0 & 0 \\
0 & 0 & 1 & \Delta t & 0 & 0 \\
0 & 0 & 0 & 1 & 0 & 0 \\
0 & 0 & 0 & 0 & 1 & \Delta t \\
0 & 0 & 0 & 0 & 0 & 1
\end{bmatrix},
\label{eq:F_matrix}
\end{equation}
where $\Delta t = 33.3$ ms is the frame interval. The process noise covariance is initialized as
\begin{equation}
\mathbf{Q} = \sigma_q^2 \mathbf{I}_6,
\label{eq:Q_matrix}
\end{equation}
with $\sigma_q^2 = 0.01$, which was selected empirically to absorb small deviations from ideal constant-velocity motion.

\subsubsection{Observation Models}

AuraBeam operates in two observation modes depending on the detector confidence.

\paragraph{Mode 1: Full observation.}
When the detector confidence is sufficiently high, the state is updated using image-plane coordinates and depth. The observation model is

\begin{equation}
\mathbf{H}_{\mathrm{full}} =
\begin{bmatrix}
1 & 0 & 0 & 0 & 0 & 0 \\
0 & 0 & 1 & 0 & 0 & 0 \\
0 & 0 & 0 & 0 & 1 & 0
\end{bmatrix},
\label{eq:H_full}
\end{equation}
with measurement vector
\begin{equation}
\mathbf{y}_k^{\mathrm{full}} = [c_x, c_y, Z(t)]^\top.
\label{eq:y_full}
\end{equation}

The corresponding measurement noise covariance is
\begin{equation}
\mathbf{R}_{\mathrm{full}} = \mathrm{diag}(2.5, 2.5, 0.5),
\label{eq:R_full}
\end{equation}
where the first two terms are measured in px$^2$ and the third term in m$^2$.

\paragraph{Mode 2: Occlusion-aware reduced observation.}
When the fused detector confidence falls below a threshold $\tau_{\mathrm{conf}}$, the image-plane observation is treated as unreliable and the filter updates only the depth component. The observation model becomes

\begin{equation}
\mathbf{H}_{\mathrm{occ}} =
\begin{bmatrix}
0 & 0 & 0 & 0 & 1 & 0
\end{bmatrix},
\label{eq:H_occ}
\end{equation}
with measurement vector
\begin{equation}
\mathbf{y}_k^{\mathrm{occ}} = [Z(t)].
\label{eq:y_occ}
\end{equation}

In this mode, the filter relies on the internal motion model to propagate $(c_x, c_y)$ while continuing to anchor the state using the depth prior. To reflect the greater uncertainty under degraded observation, the measurement covariance is inflated as
\begin{equation}
\mathbf{R}_{\mathrm{occ}} = 10 \mathbf{R}_{\mathrm{full}}.
\label{eq:R_occ}
\end{equation}

Conceptually, this mode switch allows the estimator to preserve short-term localization continuity without trusting unstable image measurements during glare-induced detector collapse.

\subsubsection{Recursive Estimation}

The prediction step is given by

\begin{equation}
\hat{\mathbf{x}}_{k}^{-} = \mathbf{F}\hat{\mathbf{x}}_{k-1},
\label{eq:kf_pred_x}
\end{equation}

\begin{equation}
\mathbf{P}_{k}^{-} = \mathbf{F}\mathbf{P}_{k-1}\mathbf{F}^{\top} + \mathbf{Q},
\label{eq:kf_pred_p}
\end{equation}
where $\hat{\mathbf{x}}_{k}^{-}$ and $\mathbf{P}_{k}^{-}$ denote the prior state estimate and covariance.
The Kalman gain is then computed as
\begin{equation}
\mathbf{K}_k = \mathbf{P}_{k}^{-}\mathbf{H}^{\top} \left(\mathbf{H}\mathbf{P}_{k}^{-}\mathbf{H}^{\top} + \mathbf{R}\right)^{-1},
\label{eq:kf_gain}
\end{equation}
followed by the update equations
\begin{equation}
\hat{\mathbf{x}}_k = \hat{\mathbf{x}}_{k}^{-} + \mathbf{K}_k\left(\mathbf{y}_k - \mathbf{H}\hat{\mathbf{x}}_{k}^{-}\right),
\label{eq:kf_update_x}
\end{equation}

\begin{equation}
\mathbf{P}_k = \left(\mathbf{I} - \mathbf{K}_k\mathbf{H}\right)\mathbf{P}_{k}^{-},
\label{eq:kf_update_p}
\end{equation}
where $(\mathbf{H}, \mathbf{R}, \mathbf{y}_k)$ are selected according to the current observation mode.

\subsubsection{Adaptive Detect-then-Track Procedure}

Algorithm~\ref{alg:main} summarizes the complete inference loop. At each frame, the system first produces a fused image-space observation from the detector ensemble, then computes the Pseudo-Radar depth, predicts the next target state, selects the observation mode based on confidence, updates the state estimate, and finally maps the filtered localization to the LED control grid.

\begin{algorithm}[t]
\caption{AuraBeam detect-then-track loop with adaptive observation switching}
\label{alg:main}
\begin{algorithmic}[1]
\Require Video stream $\{I_t\}$, thresholds $\tau_{\mathrm{conf}}$ and $\tau_{\mathrm{NMS}}$, approach velocity $v_{\mathrm{app}}$
\Ensure LED suppression grid $\mathbf{G}_t$ for each frame
\State Initialize $\hat{\mathbf{x}}_0$, $\mathbf{P}_0$, $\mathbf{Q}$, $\mathbf{R}_{\mathrm{full}}$, and $\mathbf{R}_{\mathrm{occ}}$
\State Estimate initial range $Z_0$ from the first valid target acquisition
\For{each frame $t = 1,2,\ldots$}
    \State Run $\mathcal{M}_1$ and $\mathcal{M}_2$ on $I_t$ to obtain $\mathcal{D}_1$ and $\mathcal{D}_2$
    \State Fuse detections using WeightedNMS to obtain $\mathbf{d}_t$
    \State Compute depth prior $Z(t) = Z_0 - v_{\mathrm{app}} t$
    \State Predict $\hat{\mathbf{x}}^-_t$ and $\mathbf{P}^-_t$ using Eqs.~\eqref{eq:kf_pred_x} and \eqref{eq:kf_pred_p}
    \If{$\hat{s}_t \geq \tau_{\mathrm{conf}}$}
        \State Set $\mathbf{H} \gets \mathbf{H}_{\mathrm{full}}$, $\mathbf{R} \gets \mathbf{R}_{\mathrm{full}}$, and $\mathbf{y}_t \gets \mathbf{y}_t^{\mathrm{full}}$
    \Else
        \State Set $\mathbf{H} \gets \mathbf{H}_{\mathrm{occ}}$, $\mathbf{R} \gets \mathbf{R}_{\mathrm{occ}}$, and $\mathbf{y}_t \gets \mathbf{y}_t^{\mathrm{occ}}$
    \EndIf
    \State Update $\hat{\mathbf{x}}_t$ and $\mathbf{P}_t$ using Eqs.~\eqref{eq:kf_gain}--\eqref{eq:kf_update_p}
    \State Map $(\hat{c}_x, \hat{c}_y, \hat{Z})$ to the LED grid to obtain $\mathbf{G}_t$
    \State Transmit $\mathbf{G}_t$ to the actuation layer
\EndFor
\end{algorithmic}
\end{algorithm}

\subsection{Module D: Zone Mapping and Anti-Flicker Scheduling}
\label{subsec:zone}

\subsubsection{Pixel-to-Grid Projection}

The filtered centroid $(\hat{c}_x, \hat{c}_y)$ is projected onto the $8\times8$ LED grid as

\begin{equation}
g_{\mathrm{col}} = \left\lfloor \frac{\hat{c}_x}{W_{\mathrm{frame}}} \cdot 8 \right\rfloor, \qquad
g_{\mathrm{row}} = \left\lfloor \frac{\hat{c}_y}{H_{\mathrm{frame}}} \cdot 8 \right\rfloor,
\label{eq:zone_map}
\end{equation}
where $W_{\mathrm{frame}}$ and $H_{\mathrm{frame}}$ are the image width and height. A square suppression region $\mathcal{B}_t$ centered at $(g_{\mathrm{col}}, g_{\mathrm{row}})$ is then formed using a half-width parameter $\delta$ measured in LED cells.

\subsubsection{Depth-Dependent Suppression Policy}

The suppression mode is modulated according to the estimated target range $\hat{Z}$. When the target is far away, no suppression is applied. At intermediate range, a dimmed suppression mode is used. At close range, the selected LED region is fully turned off. The thresholds used in this work are summarized in Table~\ref{tab:z_thresh}.

\begin{table}[t]
\centering
\caption{Depth-dependent glare suppression policy.}
\label{tab:z_thresh}
\begin{tabular}{lll}
\hline
\textbf{Estimated range} & \textbf{LED state} & \textbf{Duty cycle} \\
\hline
$\hat{Z} \geq 60$ m & No suppression & 100\% \\
$30 \leq \hat{Z} < 60$ m & Partial suppression & 40\% \\
$\hat{Z} < 30$ m & Full suppression & 0\% \\
\hline
\end{tabular}
\end{table}

\subsubsection{Hysteresis-Based Anti-Flicker Control}

Direct frame-by-frame switching of LED cells can produce visible flicker when the estimated suppression region oscillates near cell boundaries. To stabilize actuation, we apply a temporal hold mechanism such that a cell remains suppressed if it belonged to the recent suppression history. Formally,

\begin{equation}
\mathbf{G}_t[r,c] = \mathrm{OFF}
\quad \text{if} \quad
\exists \tau \in [t-T_h, t] \ \text{such that} \ (r,c) \in \mathcal{B}_{\tau},
\label{eq:hysteresis}
\end{equation}
where $T_h$ is the hold time expressed in frames. In all experiments, we set $T_h = 5$, corresponding to approximately 167 ms at the operating frame rate. This hysteresis mechanism reduces visible ghosting and prevents rapid toggling of adjacent cells during short-term localization jitter.

\subsection{Module E: Hardware Actuation Layer}
\label{subsec:hw}

The final suppression pattern is transmitted to an Arduino Uno, which scans the LED matrix using row-column multiplexing at 200 Hz. This refresh rate is well above the critical flicker fusion threshold of the human visual system and therefore ensures stable visual output. Communication from the host computer uses a 115200-baud serial link with the packet format

\begin{verbatim}
BOX:<col_start>:<col_end>:<row_start>:<row_end>
\end{verbatim}

The serial receiver is implemented using an interrupt-driven ring buffer so that packet decoding does not block the timer interrupt responsible for LED scanning. As a result, the actuation layer remains stable even when the upstream perception pipeline exhibits variable inference latency.

\subsection{Implementation Summary}
\label{subsec:implementation_summary}

In summary, AuraBeam combines a dual-model nighttime detector, a lightweight depth surrogate, an adaptive 3-D state estimator, and a discrete actuation controller within a single closed-loop HIL framework. The overall design is intended to evaluate a specific failure mode of camera-centric ADB, namely the loss of stable suppression under headlight saturation, while remaining sufficiently simple and reproducible for controlled experimentation.

\section{Evaluation Protocol}
\label{sec:eval}
%% ==============================================================
This section defines the protocol used to evaluate AuraBeam as a
closed-loop perception-and-actuation system. The design of the protocol
follows the three system-level claims of the paper: \emph{(i)}
ensemble detection should improve nighttime robustness under glare and
blooming, \emph{(ii)} adaptive 3-D tracking should preserve
localization during temporary visual degradation, and \emph{(iii)}
these gains should be reflected at the level of LED-grid glare
suppression rather than only in image-space detection scores.

\subsection{Implementation Details}
\label{subsec:eval_impl}

All experiments use the complete AuraBeam pipeline described in
Section~\ref{sec:method}. The perception front-end consists of two
independently trained detectors, namely YOLOv5 ($\mathcal{M}_1$) and
YOLOv8 ($\mathcal{M}_2$), whose outputs are fused by Weighted NMS with
$\tau_{\mathrm{NMS}} = 0.5$. As described in
Section~\ref{subsubsec:training}, $\mathcal{M}_1$ is trained on a
5{,}500-image headlight-focused nighttime dataset, whereas
$\mathcal{M}_2$ is fine-tuned on a second in-house nighttime dataset of
6{,}000 manually annotated images. Both detectors use an input
resolution of $1024\times1024$, batch size 16, and a maximum of 100
epochs; the ensemble weights $w_1$ and $w_2$ are selected on the
validation split and then fixed for all downstream experiments.

The tracking layer uses the six-state Kalman filter in
Section~\ref{subsec:kf_fusion}, together with the Pseudo-Radar depth
surrogate in Section~\ref{subsec:radar}. During inference, the
observation model switches according to detector confidence
$\tau_{\mathrm{conf}}$, allowing the filter to fall back to depth-only
updates during temporary visual collapse. The actuation layer maps the
filtered target state onto an $8\times8$ LED grid and applies the
hysteresis controller of Section~\ref{subsec:zone} with hold-time
$T_h = 5$ frames before transmitting the final suppression pattern to
the Arduino-based HIL platform. Unless otherwise stated, all detector,
tracking, and control parameters are fixed before test-time evaluation
and are not retuned for individual scenarios.

\subsection{Evaluation Metrics}
\label{subsec:eval_metrics}

The evaluation uses three levels of metrics. First, detector-level
metrics quantify perception quality in image space. Second,
trajectory-level metrics quantify temporal stability of the estimated
target path. Third, control-level metrics quantify whether the final LED
suppression pattern correctly covers the opposing headlight region.
Because ADB is ultimately an actuation problem, the control-level
metrics are treated as primary and the image-space metrics as
supporting evidence.

\subsubsection{Primary Control-Level Metric}
% Viết lại theo tổ chức như thế này
% removed draft placeholder
% thông số kỹ thuật, máy để train, train như thế nào,...
% removed draft placeholder
% các thông số đánh giá$
% removed draft placeholder
% xây dựng data set và chọn base line, các kịch bản thực nghiệm 
% removed draft placeholder
% thực nghiệm
% So sánh, ablation studies, phân tích,...

% đây là bản cũ
% removed draft placeholder
\label{subsec:gssr}

Conventional metrics such as mAP, Precision, and Recall characterize
detection accuracy in pixel space. For a matrix headlight controller,
however, the terminal output is a set of discrete LED cells: a small
pixel-space error may map to the wrong LED cell being suppressed,
resulting in residual glare or unnecessary darkening. A metric defined
directly over the LED grid is therefore required. The GSSR metric is
proposed to address this gap:

\begin{equation}
\mathrm{GSSR} = \frac{1}{N} \sum_{t=1}^{N}
\mathbf{1}\!\left[
  \mathrm{IoU}_{\mathrm{grid}}\!\left(\hat{\mathcal{B}}_t,\,
  \mathcal{B}^{*}_t\right) \geq 0.5
\right],
\label{eq:gssr}
\end{equation}

\noindent where $\hat{\mathcal{B}}_t$ is the set of LED cells predicted
by the system to require suppression at time $t$, and
$\mathcal{B}^{*}_t$ is the ground-truth set of cells corresponding to
the actual headlight footprint of the oncoming vehicle on the
$8 \times 8$ grid. Grid-level IoU is defined as:

\begin{equation}
\mathrm{IoU}_{\mathrm{grid}}(A,B) =
\frac{|A \cap B|}{|A \cup B|}.
\label{eq:grid_iou}
\end{equation}

The IoU threshold of 0.5 ensures that more than half of the glare
region is covered, satisfying the minimum criterion for glare reduction
specified in UN ECE R123 Article 6.2.

\subsection{Supplementary Metrics}
\label{subsec:metrics}

\noindent\textbf{False Darkening Rate (FDR):}
The proportion of frames in which the system suppresses at least one
LED cell when no opposing vehicle is present:

\begin{equation}
\mathrm{FDR} = \frac{1}{N_0}
\sum_{t:\, \mathcal{B}^{*}_t = \emptyset}
\mathbf{1}\!\left[\hat{\mathcal{B}}_t \neq \emptyset\right].
\label{eq:fdr}
\end{equation}

This metric quantifies unnecessary darkening, which degrades the host
driver's forward visibility. ~\cite{Lau2025FAR}

\medskip
\noindent\textbf{Missed Glare Rate (MGR):}
The proportion of frames in which the system fails to suppress any LED
cell when an opposing vehicle is present:

\begin{equation}
\mathrm{MGR} = \frac{1}{N_1}
\sum_{t:\, \mathcal{B}^{*}_t \neq \emptyset}
\mathbf{1}\!\left[\hat{\mathcal{B}}_t = \emptyset\right].
\label{eq:mgr}
\end{equation}

MGR is the primary safety metric, representing events in which the
system fails to prevent glare exposure.

\medskip
\noindent\textbf{Trajectory RMSE:}
\begin{equation}
\mathrm{RMSE} = \sqrt{
  \frac{1}{N} \sum_{t=1}^{N}
  \left\|
    \begin{bmatrix}\hat{c}_{x,t} \\ \hat{c}_{y,t}\end{bmatrix}
    -
    \begin{bmatrix}c^{*}_{x,t} \\ c^{*}_{y,t}\end{bmatrix}
  \right\|_2^2
},
\label{eq:rmse}
\end{equation}

\noindent where $(c^{*}_{x,t}, c^{*}_{y,t})$ denotes the
ground-truth centroid coordinates. ~\cite{Bishop2006}

\medskip
\noindent\textbf{Jitter Reduction (JR\%):}
\begin{equation}
\mathrm{JR} = \left(1 -
  \frac{\sigma_{\hat{c}}^2}{\sigma_{d}^2}
\right) \times 100\%,
\label{eq:jitter}
\end{equation}

\noindent where $\sigma_{\hat{c}}^2$ and $\sigma_{d}^2$ are the
centroid position variances of the Kalman Filter output and the raw
YOLO output, respectively.

\medskip
In the results section, mAP@0.5, Precision, and Recall are used to
analyze detector behavior; RMSE and JR\% are used to analyze tracking
quality; and GSSR, FDR, and MGR are used to analyze the final control
outcome. This separation is intentional, because a detector can trade a
small amount of mAP for a substantial gain in Recall and still improve
LED-level glare suppression.

\subsection{Dataset, Scenarios, and Baseline Configurations}
\label{subsec:eval_data_baseline}

Evaluation is performed on a manually curated \emph{Golden Test Set}
that is reserved for final benchmarking and is not used for detector
training. Pixel-level bounding boxes are annotated in Roboflow and then
projected onto LED grid cells via Eq.~\eqref{eq:zone_map} to obtain
actuation-level ground truth. The benchmark is intentionally organized
around failure-oriented scenarios rather than generic nighttime driving,
because the objective of this paper is to evaluate robustness under
camera-centric ADB failure modes.

To separate the contribution of the main modules, the experiments use
three groups of configurations: perception variants, tracking/fusion
baselines, and control variants. The perception study compares the two
single detectors against their fused ensemble. The tracking study
compares raw detection, 2-D Kalman filtering, and the full AuraBeam
configuration. The control study compares a minimal actuation strategy
against the proposed adaptive switching and anti-flicker logic.
Table~\ref{tab:eval_configs} summarizes the principal configurations
used throughout the paper.

\begin{table}[t]
\centering
\caption{Evaluation configurations used in the experimental protocol}
\label{tab:eval_configs}
\begin{tabular}{lll}
\hline
\textbf{Group} & \textbf{Configuration} & \textbf{Purpose} \\
\hline
Perception & $\mathcal{M}_1$ only & Single-model baseline \\
Perception & $\mathcal{M}_2$ only & Complementary single-model baseline \\
Perception & Ensemble & Proposed detector fusion \\
Tracking & YOLO only & No temporal filtering \\
Tracking & YOLO + KF (2D) & Image-space smoothing baseline \\
Tracking & AuraBeam & Full 3-D fusion with switching \\
Control & Fixed observation / no hold-time & Minimal control baseline \\
Control & Adaptive switching + hold-time & Proposed control logic \\
\hline
\end{tabular}
\end{table}

This baseline structure supports the interpretation of later results.
Improvements from single-model detection to the ensemble are attributed
to perception robustness; improvements from YOLO-only to YOLO + KF (2D)
are attributed to temporal smoothing; and improvements from filtered
tracking to full AuraBeam are attributed to the combined effect of 3-D
state estimation, Pseudo-Radar depth continuity, and adaptive
observation switching during visual dropout.

\subsubsection{Golden Test Set Construction}
\label{subsec:testset}

A \emph{Golden Test Set} of 913 video frames is constructed via manual
annotation using Roboflow. Pixel-level bounding
boxes are subsequently projected onto LED grid cells via
Eq.~\eqref{eq:zone_map}. Table~\ref{tab:testset} summarizes the five
scenario groups.

Sequences are selected to represent adversarial edge cases that induce
failure in camera-only ADB systems:

\begin{itemize}
    \item \textbf{Urban Night.} Standard nighttime driving with
    unobstructed sightlines, serving as the control group for baseline
    detection performance.

    \item \textbf{Curved Roads.} Unlit mountain passes with successive
    bends and oncoming traffic, designed to stress the Kalman filter
    under nonlinear trajectories and evaluate lateral tracking error.

    \item \textbf{Rainy Night (Specular Reflection Noise).} Wet
    pavement generates specular reflections from opposing headlights,
    producing spurious light sources that challenge the vision pipeline's
    FDR suppression capability.

    \item \textbf{Fog/Snow (Scattering and Blooming).} Airborne
    particles scatter light and create diffuse bloom halos that
    obliterate vehicle geometric features. This scenario highlights the
    role of the ensemble, where $\mathcal{M}_1$ (light-source
    specialist) compensates for $\mathcal{M}_2$'s shape-feature
    degradation.

    \item \textbf{Lightning Flash (Exposure Disruption).} Sudden
    high-intensity illumination transients trigger rapid AEC adjustment,
    causing multi-frame detection dropout. This is the primary stress
    test for the 3-D Kalman filter's ability to sustain trajectory
    continuity via Pseudo-Radar during complete visual signal loss.
\end{itemize}

\begin{table}[t]
\centering
\caption{Golden Test Set scenario composition and environmental
         challenges}
\label{tab:testset}
%
    \begin{tabular}{llcl}
    \hline
    \textbf{Scenario} & \textbf{Challenge Type}
      & \textbf{Frames} & \textbf{Primary Difficulty} \\
    \hline
    Urban night      & Baseline                      & 204
      & Minimal environmental noise \\
    Curved road      & Nonlinear trajectory          & 78
      & Rapid lateral displacement \\
    Rainy night      & Specular reflection noise     & 219
      & Road reflections and lens water droplets \\
    Fog / Snow       & Strong light scattering       & 173
      & Blooming eliminates object features \\
    Lightning flash  & Abrupt illumination transient & 239
      & AEC disruption causes multi-frame dropout \\
    \hline
    \textbf{Total}   &                               & \textbf{913} & \\
    \hline
    \end{tabular}%
\end{table}

Fig.~\ref{fig:testset} illustrates representative frames from each
scenario, depicting the specular reflection, bloom scattering, and
exposure disruption conditions that constitute the primary failure
triggers for camera-only systems.

\begin{figure}[t]
\centering
\begin{subfigure}{0.3\linewidth}
    \includegraphics[width=\linewidth]{figs/night_frame.png}
    \caption*{(a)}
\end{subfigure}
\hfill
\begin{subfigure}{0.3\linewidth}
    \includegraphics[width=\linewidth]{figs/curved_frame.png}
    \caption*{(b)}
\end{subfigure}
\hfill
\begin{subfigure}{0.3\linewidth}
    \includegraphics[width=\linewidth]{figs/raining_frame1.png}
    \caption*{(c)}
\end{subfigure}

\vspace{0.3cm}

\begin{subfigure}{0.3\linewidth}
    \includegraphics[width=\linewidth]{figs/raining_frame2.png}
    \caption*{(d)}
\end{subfigure}
\hfill
\begin{subfigure}{0.3\linewidth}
    \includegraphics[width=\linewidth]{figs/snow_frame.png}
    \caption*{(e)}
\end{subfigure}
\hfill
\begin{subfigure}{0.3\linewidth}
    \includegraphics[width=\linewidth]{figs/thunder_frame1.png}
    \caption*{(f)}
\end{subfigure}

\vspace{0.3cm}

\begin{subfigure}{0.3\linewidth}
    \includegraphics[width=\linewidth]{figs/thunder_frame2.png}
    \caption*{(g)}
\end{subfigure}

\caption{Representative Golden Test Set frames.
(a)~Urban night; (b)~Curved road; (c--d)~Rainy night;
(e)~Fog/Snow; (f--g)~Lightning flash.}
\label{fig:testset}
\end{figure}

\subsection{Reporting Protocol}
\label{subsec:eval_protocol}

All results reported in Section~\ref{sec:results} are computed on the
same Golden Test Set with identical post-processing and actuation
settings. Scenario-wise reporting is retained because the dominant
failure mechanisms differ substantially across rain, scattering,
curvature, and abrupt exposure disruption. No scenario-specific
threshold tuning is performed at test time. This fixed-protocol
evaluation is intended to expose genuine robustness rather than gains
obtained from per-scenario optimization.

\begin{figure}[t]
\centering
\begin{subfigure}{0.3\linewidth}
    \includegraphics[width=\linewidth]{figs/night_frame.png}
    \caption*{(a)}
\end{subfigure}
\hfill
\begin{subfigure}{0.3\linewidth}
    \includegraphics[width=\linewidth]{figs/curved_frame.png}
    \caption*{(b)}
\end{subfigure}
\hfill
\begin{subfigure}{0.3\linewidth}
    \includegraphics[width=\linewidth]{figs/raining_frame1.png}
    \caption*{(c)}
\end{subfigure}

\vspace{0.3cm}

\begin{subfigure}{0.3\linewidth}
    \includegraphics[width=\linewidth]{figs/raining_frame2.png}
    \caption*{(d)}
\end{subfigure}
\hfill
\begin{subfigure}{0.3\linewidth}
    \includegraphics[width=\linewidth]{figs/snow_frame.png}
    \caption*{(e)}
\end{subfigure}
\hfill
\begin{subfigure}{0.3\linewidth}
    \includegraphics[width=\linewidth]{figs/thunder_frame1.png}
    \caption*{(f)}
\end{subfigure}

\vspace{0.3cm}

\begin{subfigure}{0.3\linewidth}
    \includegraphics[width=\linewidth]{figs/thunder_frame2.png}
    \caption*{(g)}
\end{subfigure}

\caption{Representative Golden Test Set frames.
(a)~Urban night; (b)~Curved road; (c--d)~Rainy night;
(e)~Fog/Snow; (f--g)~Lightning flash.}
\label{fig:testset}
\end{figure}

%% ==============================================================
\section{Experimental Results}
\label{sec:results}
%% ==============================================================

\subsection{Detection Performance (Module A)}
\label{subsec:res_det}

Table~\ref{tab:det_all} presents detection results for the three
model configurations---$\mathcal{M}_1$, $\mathcal{M}_2$, and
Ensemble---across all five Golden Test Set scenarios.

\begin{table}[t]
\centering
\caption{Detection performance comparison: single models vs.\ Ensemble
         across scenarios}
\label{tab:det_all}
%
\begin{tabular}{llccc}
\hline
\textbf{Scenario} & \textbf{Configuration}
  & \textbf{mAP@0.5} & \textbf{Precision} & \textbf{Recall} \\
\hline
\multirow{3}{*}{Urban night}
 & $\mathcal{M}_1$   & 53.12\% & 69.66\% & 52.57\% \\
 & $\mathcal{M}_2$   & 48.52\% & 68.82\% & 45.61\% \\
 & \textbf{Ensemble} & \textbf{48.74\%} & \textbf{53.57\%}
                     & \textbf{56.93\%} \\
\hline
\multirow{3}{*}{Curved road}
 & $\mathcal{M}_1$   & 16.16\% & 55.32\% & 20.97\% \\
 & $\mathcal{M}_2$   & 45.08\% & 73.81\% & 50.00\% \\
 & \textbf{Ensemble} & \textbf{59.98\%} & \textbf{50.32\%}
                     & \textbf{62.90\%} \\
\hline
\multirow{3}{*}{Rainy night}
 & $\mathcal{M}_1$   & 16.70\% & 31.97\% & 19.20\% \\
 & $\mathcal{M}_2$   & 32.42\% & 64.81\% & 27.62\% \\
 & \textbf{Ensemble} & \textbf{29.76\%} & \textbf{47.91\%}
                     & \textbf{38.82\%} \\
\hline
\multirow{3}{*}{Fog / Snow}
 & $\mathcal{M}_1$   &  9.92\% & 12.61\% & 42.25\% \\
 & $\mathcal{M}_2$   & 18.14\% & 30.31\% & 28.01\% \\
 & \textbf{Ensemble} & \textbf{24.60\%} & \textbf{41.30\%}
                     & \textbf{44.88\%} \\
\hline
\multirow{3}{*}{Lightning flash}
 & $\mathcal{M}_1$   & 39.76\% & 39.27\% & 54.88\% \\
 & $\mathcal{M}_2$   &  5.36\% & 13.85\% &  8.37\% \\
 & \textbf{Ensemble} & \textbf{49.51\%} & \textbf{31.45\%}
                     & \textbf{68.40\%} \\
\hline
\end{tabular}%
\end{table}

\subsubsection*{The mAP--Recall Tradeoff: An IoU Drift Analysis}
\label{subsubsec:tradeoff}

A notable pattern in Table~\ref{tab:det_all} is that the Ensemble
mAP@0.5 does \emph{not} strictly dominate the best single model in
every scenario (e.g., Urban night: 48.74\% vs.\ 53.12\% for
$\mathcal{M}_1$). This outcome is a \emph{mathematically predictable
consequence} of Weighted NMS and is analyzed as follows.

When two overlapping boxes $(b_a, s_a)$ from $\mathcal{M}_1$ and
$(b_b, s_b)$ from $\mathcal{M}_2$ satisfy
$\mathrm{IoU}(b_a, b_b) \geq \tau_{\mathrm{NMS}}$, the fused box
$\hat{b}$ is the confidence-weighted centroid
(Eq.~\eqref{eq:nms_box}). Since $b_a \neq b_b$ in the general case,
$\hat{b}$ is displaced (drifted) from both source boxes and, by
extension, from the ground-truth annotation $b^*$. Formally, letting
$\Delta b = \hat{b} - b^*$ denote the residual displacement,
$\mathrm{IoU}(\hat{b}, b^*)$ decreases monotonically with
$\|\Delta b\|$. In practice, this IoU drift pushes borderline samples
below the 0.5 threshold, converting True Positives into False
Positives and thereby eroding mAP@0.5.

This tradeoff is, however, \textbf{by design}. The AuraBeam
architecture is governed by an ADAS fail-safe
philosophy~\cite{ISO26262_2018}:
\begin{itemize}[nosep]
  \item A \textbf{False Negative}---failing to detect an oncoming
    light source---causes the suppression zone to disappear, restoring
    full high-beam toward the opposing driver and creating a direct
    safety hazard.
  \item A \textbf{False Positive}---suppressing an LED cell when no
    vehicle is present---reduces the host driver's illumination
    marginally, without a safety consequence.
\end{itemize}
Under the principle of \emph{preferring over-suppression to
under-suppression}, running $\mathcal{M}_1$ and $\mathcal{M}_2$ in
parallel yields substantial Recall gains across adversarial scenarios:
from 8.37\% ($\mathcal{M}_2$ alone) to 68.40\% (Ensemble) in the
lightning scenario, and from 19.20\% to 38.82\% in the rainy night
scenario. In this context, \textbf{maximizing Recall carries a more
direct safety implication than optimizing mAP}, and the proposed GSSR
metric (Section~\ref{subsec:gssr}) is the appropriate evaluation
standard for LED-level actuation fidelity.

Four scenario-level observations follow. \textbf{(i)}~In Urban night
(baseline), $\mathcal{M}_1$ achieves peak mAP (53.12\%) under stable
illumination; the Ensemble improves Recall by 4.4 percentage points
(56.93\%) at an acceptable mAP cost. \textbf{(ii)}~In the Curved road
scenario, feature complementarity between the two models produces the
largest mAP gain ($+$14.9 points over $\mathcal{M}_2$), confirming
ensemble efficacy under nonlinear trajectories. \textbf{(iii)}~In the
Fog/Snow scenario, severe blooming reduces $\mathcal{M}_1$ Precision
to 12.61\%; the Ensemble recovers to 41.30\%, demonstrating
multi-model resilience. \textbf{(iv)}~The lightning scenario most
directly validates the dual-model architecture: $\mathcal{M}_2$
nearly collapses (mAP 5.36\%) while the Ensemble sustains 68.40\%
Recall, constituting direct empirical evidence that a single-model
architecture is insufficient for this problem.

\subsection{Tracking and Fusion Performance (Module C)}
\label{subsec:res_kf}

Table~\ref{tab:tracking_scenarios} compares trajectory RMSE and
jitter reduction across the three tracker configurations.
Fig.~\ref{fig:traj_lightning} and Fig.~\ref{fig:traj_rain} illustrate
representative trajectory plots for two characteristic occlusion
scenarios.

\begin{table}[t]
\centering
\caption{Tracking performance across the five Golden Test scenarios}
\label{tab:tracking_scenarios}

\begin{tabular}{l ccc ccc}
\hline
\multirow{2}{*}{\textbf{Scenario}}
  & \multicolumn{3}{c}{\textbf{RMSE (px) $\downarrow$}}
  & \multicolumn{3}{c}{\textbf{JR\% $\uparrow$}} \\
\cmidrule(lr){2-4} \cmidrule(lr){5-7}
  & YOLO & {+KF (2D)} & {AuraBeam}
  & YOLO & {+KF (2D)} & {AuraBeam} \\
\hline
Urban night    & 36.8  & 16.5 & \textbf{16.5}
               & ---   & $-$3.4  & \textbf{$-$3.4} \\
Curved road    & 98.5  & 36.1 & \textbf{36.1}
               & ---   & 2.9     & \textbf{2.9} \\
Rainy night    & 238.1 & 80.2 & \textbf{80.2}
               & ---   & $-$10.5 & \textbf{$-$10.5} \\
Fog / Snow     & 86.9  & 25.1 & \textbf{25.1}
               & ---   & 12.7    & \textbf{12.7} \\
Lightning      & 78.7  & 11.7 & \textbf{11.7}
               & ---   & 16.3    & \textbf{16.3} \\
\hline
\textbf{Mean}  & 107.8 & 33.9 & \textbf{33.9}
               & ---   & 3.6     & \textbf{3.6} \\
\hline
\end{tabular}
\end{table}

\noindent\textbf{Note on the AuraBeam column.} The ``AuraBeam'' column
reflects the full configuration (Ensemble YOLO + 3-D KF +
Pseudo-Radar). In scenarios without complete occlusion (Urban night,
Curved road), the Pseudo-Radar contribution to planar $(x,y)$ RMSE
is negligible; its decisive role manifests along the Z-axis during
camera signal loss, as evidenced by the GSSR jump in
Table~\ref{tab:ablation}.

The 3-D Kalman filter reduces mean RMSE from 107.8\,px (YOLO only)
to 33.9\,px, a 68.6\% reduction. The highest JR\% values occur in the
lightning (16.3\%) and fog/snow (12.7\%) scenarios, confirming that
filter smoothing is most effective precisely when the vision signal is
most degraded. Negative JR\% in the urban night and rainy night
scenarios reflects a small phase lag introduced by the filter under
noise distributions that deviate from the assumed process model---an
acceptable tradeoff relative to the occlusion stability benefit.

\begin{figure}[t]
    \centering
    \includegraphics[width=0.95\linewidth]{lightning_ablation_trajectory.pdf}
    \caption{Centroid trajectory and RMSE analysis for the Lightning
             Flash scenario. The shaded region denotes the period of
             camera AEC disruption during which YOLO detection drops
             out entirely.}
    \label{fig:traj_lightning}
\end{figure}

\begin{figure}[t]
    \centering
    \includegraphics[width=0.95\linewidth]{raining_ablation_trajectory.pdf}
    \caption{Centroid trajectory and RMSE analysis for the Heavy Rain
             scenario. AuraBeam sustains stable trajectory prediction
             across 88 consecutive frames of complete detection dropout.}
    \label{fig:traj_rain}
\end{figure}

\subsection{Occlusion Resilience: GSSR Results}
\label{subsec:res_gssr}

Table~\ref{tab:gssr} presents the primary result of this paper,
comparing the three system configurations across all five Golden Test
Set scenarios on the three actuation-level metrics: GSSR, FDR, and
MGR (defined in Section~\ref{subsec:metrics}).

\begin{table}[t]
\centering
\caption{GSSR, FDR, and MGR across the five test scenarios (\%)}
\label{tab:gssr}
%
\begin{tabular}{llccc}
\hline
\textbf{Scenario} & \textbf{Configuration}
  & \textbf{GSSR}$\uparrow$ & \textbf{FDR}$\downarrow$
  & \textbf{MGR}$\downarrow$ \\
\hline
\multirow{3}{*}{Urban night}
  & YOLO only         & 94.0\% &  4.9\% & 0.0\% \\
  & YOLO + KF (2D)    & 92.0\% &  7.4\% & 0.0\% \\
  & \textbf{AuraBeam} & \textbf{96.0\%} & \textbf{4.4\%}
                      & \textbf{0.0\%} \\
\hline
\multirow{3}{*}{Curved road}
  & YOLO only         & 75.0\%  &  7.9\% & 3.9\% \\
  & YOLO + KF (2D)    & 80.3\%  &  3.9\% & 0.0\% \\
  & \textbf{AuraBeam} & \textbf{84.2\%} & \textbf{15.8\%}
                      & \textbf{0.0\%} \\
\hline
\multirow{3}{*}{Rainy night}
  & YOLO only         & 46.1\%  & 12.8\% & 40.2\% \\
  & YOLO + KF (2D)    & 82.6\%  &  3.2\% &  3.2\% \\
  & \textbf{AuraBeam} & \textbf{87.2\%} & \textbf{9.6\%}
                      & \textbf{3.2\%} \\
\hline
\multirow{3}{*}{Fog / Snow}
  & YOLO only         & 97.1\%  & 0.0\% & 0.0\% \\
  & YOLO + KF (2D)    & 100.0\% & 0.0\% & 0.0\% \\
  & \textbf{AuraBeam} & \textbf{98.8\%} & \textbf{1.2\%}
                      & \textbf{0.0\%} \\
\hline
\multirow{3}{*}{Lightning flash}
  & YOLO only         & 34.7\%  &  0.4\% & 4.6\% \\
  & YOLO + KF (2D)    & 36.4\%  &  1.3\% & 0.0\% \\
  & \textbf{AuraBeam} & \textbf{99.2\%} & \textbf{0.8\%}
                      & \textbf{0.0\%} \\
\hline
\end{tabular}%
\end{table}

Two principal findings emerge from Table~\ref{tab:gssr}.

\textbf{Breakthrough in the lightning scenario.} Under sudden
exposure disruption, the YOLO-only configuration degrades to a GSSR
of 34.7\%, and the YOLO + 2-D KF configuration improves only
marginally (36.4\%), as the 2-D filter lacks Z-axis information to
interpolate the trajectory kinematically during visual signal loss.
Integrating Pseudo-Radar depth elevates GSSR to \textbf{99.2\%}---an
absolute gain of $+$62.8 percentage points over both baselines---
confirming that 3-D tracking continuity through visual occlusion is
the decisive factor in the most dangerous ADB scenario.

\textbf{FDR tradeoff on curved roads.} AuraBeam achieves MGR = 0.0\%
but incurs an elevated FDR of 15.8\%, compared to 3.9\% for the 2-D
KF baseline. This results from the constant-velocity kinematic model's
linear prediction inertia: the filter maintains the bounding box at
the projected location even after the vehicle completes a curve.
Under the ADAS fail-safe principle---tolerating False Positives
(over-suppression) to prevent False Negatives (missed glare)---this
behavior is an intentional design property. A constant-turn-rate
kinematic model is identified as the primary mitigation in future
work.

Across the full Golden Test Set, AuraBeam achieves MGR = 0.0\% in 4
of 5 scenarios. The single exception (Rainy night, MGR = 3.2\%) still
represents a 37 percentage point reduction relative to the YOLO-only
baseline (40.2\%). These results confirm that the proposed 3-D sensor
fusion architecture satisfies the core safety requirement of the ADB
problem: \emph{no unmitigated glare exposure events} under adversarial
environmental conditions.

\subsection{Hardware Latency and Real-Time Performance}
\label{subsec:res_hw}

Table~\ref{tab:latency_all} presents the processing time budget across
the five scenarios. YOLO inference time includes preprocessing and
NMS post-processing for both parallel models, evaluated on a CPU-only
host in an offline setting.

\begin{table}[t]
\centering
\caption{Software latency and throughput analysis across five scenarios
         (offline evaluation environment)}
\label{tab:latency_all}
%
\begin{tabular}{lcccc}
\hline
\textbf{Scenario}
  & \textbf{YOLO Inference (ms)}
  & \textbf{AuraBeam Fusion (ms)}
  & \textbf{Total Latency (ms)}
  & \textbf{FPS (offline)} \\
\hline
Urban night    & 421.5 & 0.90 & 429.6 & 2.32 \\
Curved road    & 379.0 & 0.60 & 388.5 & 2.57 \\
Rainy night    & 396.6 & 0.50 & 403.2 & 2.48 \\
Fog / Snow     & 419.3 & 0.49 & 426.1 & 2.35 \\
Lightning      & 419.0 & 0.57 & 425.7 & 2.35 \\
\hline
\textbf{Mean}  & \textbf{407.1} & \textbf{0.61}
               & \textbf{414.6} & \textbf{2.41} \\
\hline
\end{tabular}%
\end{table}

The throughput bottleneck lies entirely in the unoptimized YOLO
inference block, which accounts for 98.2\% of total latency in the
offline environment. The AuraBeam fusion core (Modules B + C + D)
consumes a mean of \textbf{0.61\,ms}---0.15\% of the time
budget---with a standard deviation of $< 0.15$\,ms across all
weather conditions, confirming deterministic execution independent of
scene complexity.

For embedded deployment, it is necessary to account for the full
latency budget. Although the fusion kernel requires only 0.61\,ms, the
practical end-to-end latency is governed by YOLO inference
(${\sim}35$\,ms with TensorRT on Jetson Orin in optimal conditions)
and I/O communication. Serial UART at 115\,200 baud with a
${\sim}$30-byte packet yields a theoretical maximum of ${\sim}$14
messages/s (71\,ms/packet), constituting an additional bottleneck.
A realistic embedded deployment is projected to achieve 20--25\,FPS.
Meeting the ISO 26262 automotive real-time standard of
$\geq$30\,FPS~\cite{ISO26262_2018} requires: (1) TensorRT quantization,
(2) replacement of UART with CAN bus communication, and (3) latency
measurement on the actual embedded platform.

\subsection{Ablation Study}
\label{subsec:ablation}

Table~\ref{tab:ablation} quantifies the marginal contribution of each
module via a progressive component addition strategy, with GSSR and
JR\% averaged over the full Golden Test Set.

\begin{table}[t]
\centering
\caption{Ablation study: marginal contribution of each module
         (mean over the full Golden Test Set)}
\label{tab:ablation}
%
\begin{tabular}{lccc}
\hline
\textbf{System Configuration}
  & \textbf{GSSR} & \textbf{$\Delta$GSSR} & \textbf{JR\%} \\
\hline
(1) YOLO only (Baseline)          & 69.4\%  & ---            & 0.0\% \\
(2) $+$ 2-D Kalman Filter         & 78.3\%  & $+$8.9\%       & 3.6\% \\
\textbf{(3) $+$ Pseudo-Radar (AuraBeam)}
  & \textbf{93.1\%} & \textbf{$+$14.8\%} & \textbf{3.6\%} \\
\hline
\end{tabular}%
\end{table}

The ablation data reveal two distinct contribution levels.
Step~(1)$\to$(2): the 2-D Kalman filter improves geometric stability
in image space, yielding a GSSR gain of 8.9 percentage points and a
jitter reduction of JR\% = 3.6\%.
Step~(2)$\to$(3): Pseudo-Radar integration is the decisive improvement,
delivering the largest single GSSR gain ($+$14.8 points) without
altering JR\%---confirming that Pseudo-Radar does not contribute to
2-D trajectory smoothing but instead resolves an independent problem:
maintaining Z-axis continuity of the suppression zone during visual
occlusion. This result supports the central thesis that the ADB
problem requires a full 3-D state representation, not merely a 2-D
smoothing of the projected vision signal.

%% ==============================================================
\section{Conclusion}
\label{sec:conclusion}
%% ==============================================================

This paper presented \textbf{AuraBeam}, an end-to-end
Hardware-in-the-Loop smart matrix headlight system that resolves the
glare-induced detection loop failure inherent in all camera-centric ADB
architectures. Three contributions were made: \textbf{(i)} an Ensemble
YOLOv5/v8 pipeline with Weighted NMS that improves detection Recall
under blooming conditions; \textbf{(ii)} a 3-D Kalman Filter with
adaptive dual observation models fusing vision coordinates with
Pseudo-Radar depth, enabling trajectory continuation through complete
visual occlusion; and \textbf{(iii)} the GSSR metric, a
hardware-aware evaluation standard that captures actuation fidelity at
the LED grid level. AuraBeam achieves a GSSR of 93.1\%, a trajectory
RMSE of 33.9\,px, and a fusion core latency of 0.61\,ms on a \$23
hardware platform. In the lightning scenario---the most dangerous
failure mode---AuraBeam sustains 99.2\% GSSR while the YOLO-only
baseline collapses to 34.7\%, validating the indispensability of 3-D
sensor fusion for robust ADB control.

\paragraph{Limitations and Safety Considerations.}
\begin{itemize}[leftmargin=*]
  \item \textbf{Constant-Velocity Model.} The Pseudo-Radar assumes
    constant $v_{\mathrm{app}}$. Under moderate braking
    ($a = -0.5\,g$), depth estimation error accumulates as
    $\Delta Z(t) \approx \frac{1}{2}|a|t^2$, reaching ${\sim}$1--2\,m
    after 1--1.5\,s of camera occlusion. Future work must implement
    adaptive velocity estimation or higher-order kinematic models
    (e.g., constant-turn-rate). A formal safety analysis linking
    maximum tolerable $\Delta Z$ to ADB control margins has not been
    performed.
  \item \textbf{$8\times8$ LED Resolution.} The prototype operates
    with 64 segments ($11.25^\circ$ per segment), substantially coarser
    than production standards ($\geq1^\circ$ per SAE J3014). This directly
    inflates FDR; scaling to a production-grade matrix geometry is a
    prerequisite for regulatory acceptance.
  \item \textbf{Single-Vehicle Tracking.} The Kalman filter assumes
    a single oncoming vehicle. Multi-vehicle scenarios (highway
    overtaking) require independent KF tracks per target, increasing
    computational load and synchronization complexity.
  \item \textbf{GSSR Metric Validation.} The proposed GSSR (grid-level
    IoU $\geq 0.5$) captures grid coverage but has \emph{not been
    validated} against photometric standards (illuminance in lux per
    UN ECE R123 Article 6). Bridging GSSR to lux compliance requires
    laboratory photometric measurement---a critical gap for genuine
    safety claims.
\end{itemize}

\paragraph{Future Work and Production Pathway.}
\begin{enumerate}[label=\textbf{(\roman*)},leftmargin=*]
  \item \textbf{Advanced Motion Models.} Replace the constant-velocity
    assumption with constant-turn-rate (CTR) or Wiener process
    acceleration models to handle real-world vehicle maneuvers.
    Quantify GSSR on curved-road scenarios where FDR currently reaches
    15.8\%.
  \item \textbf{Real Sensor Fusion.} Integrate a 77--79\,GHz mmWave
    Radar in place of the Pseudo-Radar simulation; benchmark depth
    accuracy against LiDAR ground truth, targeting $|\Delta Z| < 0.5$\,m
    under emergency braking.
  \item \textbf{Full Embedded Deployment.} Port to NVIDIA Jetson Orin
    Nano with TensorRT optimization; measure end-to-end latency from
    camera capture to LED actuation under real-time OS constraints.
    Validate $\geq$30\,FPS compliance per ISO 26262~\cite{ISO26262_2018}.
  \item \textbf{Production LED Hardware.} Scale from the $8\times8$
    prototype (64 segments) to a production-grade matrix
    ($\geq$512 segments); re-evaluate GSSR and FDR scaling laws under
    finer spatial granularity.
  \item \textbf{Photometric Validation.} Establish a goniophotometric
    bench with calibrated luminance meters; correlate the proposed GSSR
    metric with actual illuminance reduction (lux) per UN ECE R123
    Article 6.2. This step is mandatory for regulatory compliance claims.
  \item \textbf{Functional Safety Analysis.} Conduct ISO 26262
    Functional Safety Assessment (ASIL rating, FMEA). Characterize
    failure modes (Pseudo-Radar corruption, Kalman filter divergence)
    and implement fail-safe fallback strategies (e.g., conservative
    full-beam cutoff).
\end{enumerate}


% \backmatter

% \bmhead{Supplementary information}

% If your article has accompanying supplementary file/s please state so here. 

% Authors reporting data from electrophoretic gels and blots should supply the full unprocessed scans for key as part of their Supplementary information. This may be requested by the editorial team/s if it is missing.

% Please refer to Journal-level guidance for any specific requirements.

\bmhead{Acknowledgements}

Acknowledgements are not compulsory. Where included they should be brief. Grant or contribution numbers may be acknowledged.

Please refer to Journal-level guidance for any specific requirements.

\section*{Declarations}



% \begin{appendices}

% \section{Section title of first appendix}\label{secA1}

% An appendix contains supplementary information that is not an essential part of the text itself but which may be helpful in providing a more comprehensive understanding of the research problem or it is information that is too cumbersome to be included in the body of the paper.

% %%=============================================%%
% %% For submissions to Nature Portfolio Journals %%
% %% please use the heading ``Extended Data''.   %%
% %%=============================================%%

% %%=============================================================%%
% %% Sample for another appendix section			       %%
% %%=============================================================%%

% %% \section{Example of another appendix section}\label{secA2}%
% %% Appendices may be used for helpful, supporting or essential material that would otherwise 
% %% clutter, break up or be distracting to the text. Appendices can consist of sections, figures, 
% %% tables and equations etc.

% \end{appendices}

%%===========================================================================================%%
%% If you are submitting to one of the Nature Portfolio journals, using the eJP submission   %%
%% system, please include the references within the manuscript file itself. You may do this  %%
%% by copying the reference list from your .bbl file, paste it into the main manuscript .tex %%
%% file, and delete the associated \verb+\bibliography+ commands.                            %%
%%===========================================================================================%%

\bibliography{sn-bibliography}% common bib file
%% if required, the content of .bbl file can be included here once bbl is generated
%%\input sn-article.bbl

\end{document}
